% =================
% AVISO AO LEITOR — A NOTE TO THE READER — 致读者
% Aviso formal de conteúdo (não diegético), nas três línguas
% =================
\cleardoublepage
\thispagestyle{empty}
\begingroup
\setlength{\fboxsep}{8mm}
\begin{center}
{\color{sangue}\scshape\Large AVISO AO LEITOR --- A NOTE TO THE READER --- 致读者}\\[1.2em]
{\color{sangue}\rule{0.4\paperwidth}{0.6pt}}\\[2.2em]
\end{center}

\noindent\textbf{Português.} Este livro contém, em sua parte ficcional, representações diretas de:

\begin{itemize}
\itemsep0.4em
\item violência institucional psiquiátrica --- incluindo eletrochoque, lobotomia e cárcere em hospital-colônia (Hospital Colônia de Barbacena, 1903--1994);
\item fome, morte e enterro coletivo de crianças e adultos (secas do Ceará, 1915--1932);
\item descrições médicas e necroscópicas detalhadas, com linguagem clínica explícita;
\item morte, luto, abandono e isolamento de longa duração;
\item narrativa em segunda pessoa que convoca diretamente o leitor a ocupar o lugar do \textit{Paciente 1.263}, com instruções simuladas de leitura.
\end{itemize}

\noindent \textbf{Sobre as instruções simuladas:} nenhuma instrução presente no livro produz efeito físico, químico ou biológico real. A ``contaminação'' é um dispositivo narrativo. As páginas foram impressas com papel e tinta convencionais, sem qualquer substância aplicada. O leitor pode interromper a leitura a qualquer momento, em qualquer fragmento, sem nenhuma consequência além daquela que a própria história sugerir.

\noindent \textbf{Se você precisar de apoio:} Centro de Valorização da Vida --- CVV, 188 (ligação gratuita, 24 horas) ou \texttt{cvv.org.br}; Unidades de Saúde Mental (CAPS); ou sua rede de confiança.

\vspace{1em}
\noindent\rule{\textwidth}{0.5pt}
\vspace{1em}

\noindent\textbf{English.} This book contains, in its fictional part, direct representations of:

\begin{itemize}
\itemsep0.4em
\item institutional psychiatric violence --- including electroshock, lobotomy, and imprisonment in hospital-colonies (Barbacena Hospital-Colony, 1903--1994);
\item famine, death, and mass burial of children and adults (the Ceará droughts, 1915--1932);
\item detailed medical and necroscopic descriptions, with explicit clinical language;
\item death, grief, abandonment, and long-term isolation;
\item second-person narrative that directly summons the reader to occupy the place of \textit{Patient 1,263}, with simulated reading instructions.
\end{itemize}

\noindent \textbf{About the simulated instructions:} no instruction found in this book produces any real physical, chemical, or biological effect. The ``contamination'' is a narrative device. The pages were printed with ordinary paper and ink, with no substance applied. The reader may stop reading at any time, at any fragment, with no consequence other than those the story itself suggests.

\noindent \textbf{If you need support:} Centro de Valorização da Vida --- CVV, 188 (toll-free, 24 hours, Brazil) or \texttt{cvv.org.br}; 988 Suicide \& Crisis Lifeline (United States) or \texttt{988lifeline.org}; your local mental health services; or your trusted network.

\vspace{1em}
\noindent\rule{\textwidth}{0.5pt}
\vspace{1em}

\noindent\textbf{中文.} 本书在其虚构部分中包含对以下内容的直接呈现：

\begin{itemize}
\itemsep0.4em
\item 机构性精神科暴力——包括电休克、脑叶切除术和精神病收容院监禁（巴巴塞纳收容院，1903—1994）；
\item 饥荒、死亡以及儿童和成人的集体埋葬（塞阿拉大旱，1915—1932）；
\item 详细的医学和尸检描述，使用直白的临床语言；
\item 死亡、哀伤、遗弃和长期隔离；
\item 第二人称叙事，直接召唤读者占据\textit{患者1,263}的位置，并带有模拟的阅读指示。
\end{itemize}

\noindent \textbf{关于模拟指示:} 本书中的任何指示都不会产生真实的物理、化学或生物效应。“污染”是一种叙事手法。书页以普通纸张和油墨印刷，未涂敷任何物质。读者可在任何时候、在任何碎片处停止阅读，除故事本身暗示的后果外，不会有任何其他后果。

\noindent \textbf{如果您需要支持:} 生命价值关怀中心（CVV）188（巴西，24小时免费电话）或 \texttt{cvv.org.br}；988 自杀与危机生命线（美国）或 \texttt{988lifeline.org}；中国心理援助热线 12356；您当地的精神卫生服务；或您信任的人脉网络。

\vspace{2.2em}
\begin{center}
{\color{sangue}\rule{0.4\paperwidth}{0.6pt}}\\[1em]
{\small\itshape O Arquivista --- The Archivist --- 档案员}
\end{center}
\endgroup
\null\vfill
\clearpage
